\chapter{实验设计与结果分析}
\label{cha:experiments}

本章介绍实验设计和实验结果分析。首先介绍仿真环境、对比方法和评估指标，然后呈现实验结果并进行分析。

\section{仿真环境与参数}
\label{sec:simulation-setup}

实验区域为 $700$m $\times$ 700m，系统包含 $N=5$ 架 UAV、$M=100$ 个 UE 和 $1$ 个 MBS。每个 episode 包含 $T=1000$ 个 time slots，每个 time slot 时长为 $\tau = 1$s。UAV 飞行高度 $H=100$m，最大速度 $v_{max}=15$m/s，覆盖半径 $R_c=100$m，感知范围 $460$m，最小 UAV 间距 $200$m。系统包含 25 类服务和 50 类内容文件，服务 deadline 在 $[0.65, 2.10]$s 范围内生成。MBS 固定算力 $f_{mbs} = 200 \times 10^9$ cycles/s。UAV 计算能力在 $[40, 90] \times 10^9$ cycles/s 之间随机生成。

UE 电池容量 $B_{max}=500$J，临界阈值 $B_{low}=50$J，待机功耗 $P_{static}=0.01$W。飞行功率 $P_{move}=300$W，悬停功率 $P_{hover}=150$W。WPT 发射功率 $P_{wpt}=50$W，采集增益 $G=5\times10^5$，采集效率 $\eta=0.8$。

所有 MARL 策略使用 PyTorch 实现，MAPPO 采用 2 层 MLP（hidden dim=128），attention 模块隐层维度为 64（ATTN\_HIDDEN\_DIM=64），Adam 优化器，learning rate$=3\times10^{-4}$，PPO clip $\epsilon=0.2$，discount $\gamma=0.99$，GAE $\lambda=0.95$。

正式实验采用 3 个 training seeds（42, 84, 126）和 10 个 workload seeds（42-420），每个 workload seed 运行 6 个 episodes。统计单元为 (training seed, workload seed) 的 episode mean，每个策略共有 $N=30$ 个统计样本。本文报告 mean、std、95\% CI、paired delta、paired t-test 和 Wilcoxon 检验。

\section{对比方法}
\label{sec:comparison-methods}

\subsection{主实验策略组合}

主实验比较五组策略组合，如表 \ref{tab:main-strategies} 所示。

\begin{table}[!htbp]
    \centering
    \bicaption{主实验策略组合}{Main experiment strategy combinations}
    \label{tab:main-strategies}
    \begin{tabular}{|l|c|c|c|}
        \hline
        组合方案 & 轨迹控制 & 卸载策略 & 角色 \\
        \hline
        uncoordinated\_greedy\_\_heuristic & uncoordinated greedy & heuristic & 基础参考 \\
        attention\_mappo\_\_heuristic & attention-MAPPO & heuristic & 上层贡献 \\
        uncoordinated\_greedy\_\_lower\_mappo & uncoordinated greedy & lower MAPPO & 下层贡献 \\
        attention\_mappo\_\_lower\_mappo & attention-MAPPO & lower MAPPO & 分别训练后组合 \\
        full\_hierarchical\_marl & attention-MAPPO & lower MAPPO & 完整双层联合训练 \\
        \hline
    \end{tabular}
\end{table}

\subsection{下层消融实验}

下层消融实验以 \texttt{lower\_full}（即 \texttt{attention\_mappo\_\_lower\_mappo}）为参考，如表 \ref{tab:ablation-strategies} 所示。

\begin{table}[!htbp]
    \centering
    \bicaption{下层消融实验配置}{Lower-layer ablation configurations}
    \label{tab:ablation-strategies}
    \begin{tabular}{|l|c|c|c|}
        \hline
        消融方法 & Attention & Quality mask & Lagrange \\
        \hline
        lower\_full & yes & yes & yes \\
        lower\_no\_mask & yes & no & yes \\
        lower\_no\_lagrange & yes & yes & no \\
        lower\_no\_attention & no & yes & yes \\
        \hline
    \end{tabular}
\end{table}

\subsection{额外 baseline}

为进一步界定 attention 机制和双层分解的适用边界，本文还补充报告上层 \texttt{vanilla\_mappo\_\_heuristic} 消融、attention 权重可视化和 \texttt{joint\_end\_to\_end\_mappo} 端到端联合 MAPPO baseline。由于补充实验当前采用单一 training seed（42）和 4 个 workload seeds（42, 84, 126, 168），统计规模小于主实验，因此主要用于机制分析和负向 baseline 说明。

\section{评估指标}
\label{sec:metrics}

本文使用以下评估指标：

\begin{itemize}
    \item \textbf{Reward ($R_t$)}：第 \ref{sec:optimization-objective} 节定义的系统级综合奖励
    \item \textbf{Latency ($L$)}：所有 UE 的总时延（含 $T_{penalty}=20$s 未服务惩罚）
    \item \textbf{Energy ($E$)}：UAV 侧总能耗（含飞行、悬停、计算、通信、WPT）
    \item \textbf{Fairness ($J$)}：基于各 UE 的 service\_coverage 计算 Jain 指数，衡量用户服务覆盖公平性
    \item \textbf{DSR}：deadline satisfaction rate，被服务且满足 deadline 的请求比例
    \item \textbf{offline\_rate ($O$)}：电池低于 $B_{low}$ 的 UE 比例
    \item \textbf{Offloading ratios}：local/cooperative/MBS 卸载比例
    \item \textbf{MBS load ratio}：卸载至 MBS 的服务请求数占已处理服务请求数的比例
\end{itemize}

\section{主实验结果}
\label{sec:main-results}

表 \ref{tab:main-results} 汇总了主实验中各策略的均值和标准差（mean $\pm$ std）。

\begin{table}[!htbp]
    \centering
    \bicaption{主实验结果（mean $\pm$ std, $N=30$）}{Main experiment results (mean $\pm$ std, $N=30$)}
    \label{tab:main-results}
    \resizebox{\textwidth}{!}{
    \begin{tabular}{l|rr|rr|rr|rr}
        \hline
        策略 & Reward & & Energy (M) & & DSR & & Fairness & \\
        \hline
        unco+heuristic & -5933.9 & $\pm$ 195.6 & 114.70 & $\pm$ 15.4 & 0.2734 & $\pm$ 0.029 & 0.7745 & $\pm$ 0.077 \\
        att+heuristic & \textbf{-2128.6} & $\pm$ 389.7 & 111.81 & $\pm$ 10.6 & \textbf{0.2879} & $\pm$ 0.032 & \textbf{0.9299} & $\pm$ 0.036 \\
        unco+lower & -5274.4 & $\pm$ 141.3 & 42.85 & $\pm$ 13.6 & 0.1891 & $\pm$ 0.029 & 0.7737 & $\pm$ 0.080 \\
        att+lower & \textbf{-1599.4} & $\pm$ 330.1 & 53.32 & $\pm$ 12.4 & 0.2170 & $\pm$ 0.035 & \textbf{0.9281} & $\pm$ 0.035 \\
        full\_hierarchical & \textbf{-1451.7} & $\pm$ 373.7 & 58.06 & $\pm$ 7.0 & 0.2083 & $\pm$ 0.019 & \textbf{0.9363} & $\pm$ 0.041 \\
        \hline
    \end{tabular}}
\end{table}

表 \ref{tab:main-results}（续）补充了 Offline\% 和卸载比例指标。

\begin{table}[!htbp]
    \centering
    \bicaption{主实验结果（续）}{Main experiment results (continued)}
    \label{tab:main-results-cont}
    \resizebox{\textwidth}{!}{
    \begin{tabular}{l|rr|rr|rr}
        \hline
        策略 & Offline\% & & Local\% & Coop\% & MBS load\% \\
        \hline
        unco+heuristic & 0.58\% & & 47.1\% & 50.1\% & 2.6\% \\
        att+heuristic & 0.10\% & & 47.2\% & 50.0\% & 2.7\% \\
        unco+lower & 0.53\% & & 52.6\% & 30.1\% & 17.1\% \\
        att+lower & 0.11\% & & 41.7\% & 44.6\% & 13.7\% \\
        full\_hierarchical & \textbf{0.00\%} & & 46.4\% & 42.4\% & 11.1\% \\
        \hline
    \end{tabular}}
\end{table}

表 \ref{tab:paired-comparison} 给出了 Full Hierarchical MARL 相对基础参考的配对比较。

\begin{table}[!htbp]
    \centering
    \bicaption{Full Hierarchical MARL 相对基础参考的配对比较}{Paired comparison of Full Hierarchical MARL vs baseline}
    \label{tab:paired-comparison}
    \begin{tabular}{l|c|c|c|c}
        \hline
        指标 & Mean Delta & 95\% CI & paired t $p$ & Effect \\
        \hline
        Reward & \textbf{+4482.2} & [4338.3, 4626.1] & $1.0\times10^{-32}$ & $\uparrow$ \\
        Energy & \textbf{-56.64M} & [-61.97M, -51.31M] & $1.7\times10^{-19}$ & $\uparrow$ \\
        Fairness & \textbf{+0.1617} & [0.1305, 0.1929] & $1.7\times10^{-11}$ & $\uparrow$ \\
        DSR & -0.0651 & [-0.0779, -0.0524] & $2.4\times10^{-11}$ & $\downarrow$ \\
        offline\_rate & -0.0058 & [-0.0099, -0.0017] & 0.0067 & $\uparrow$ \\
        Coop ratio & -0.0776 & [-0.1005, -0.0546] & $1.3\times10^{-7}$ & $\downarrow$ \\
        MBS load ratio & +0.0387 & [0.0319, 0.0454] & $1.7\times10^{-12}$ & 增加 \\
        \hline
    \end{tabular}
\end{table}

完整双层 MARL 相比基础参考，reward 提升 4482.2（$p = 1.0\times10^{-32}$），UAV 侧能耗降低 56.64M（$p = 1.7\times10^{-19}$），公平性提升 0.162（$p = 1.7\times10^{-11}$），offline rate 降至 0.00\%。DSR 是完整方法相对基础参考的明确负向结果：从 0.2734 降至 0.2083（$\Delta=-0.0651$, $p=2.4\times10^{-11}$），说明当前能耗优先配置下完整方法主要获得的是能耗、公平性和在线率收益，而非 DSR 收益。

\begin{figure}[!htbp]
    \centering
    \includegraphics[width=0.85\textwidth]{docs/figures/fig_main_comparison.png}
    \bicaption{主实验多指标对比}{Main experiment multi-metric comparison}
    \label{fig:main-comparison}
\end{figure}

\section{上层贡献分析}
\label{sec:upper-contribution}

仅替换上层轨迹策略（\texttt{att+heuristic} vs \texttt{uncoordinated\_greedy\_\_heuristic}），在保持 heuristic 卸载不变的情况下：
\begin{itemize}
    \item Reward：-5933.9 $\rightarrow$ -2128.6（$\Delta=+3805.4$, $p=3.0\times10^{-30}$）
    \item Fairness：0.7745 $\rightarrow$ 0.9299（$\Delta=+0.1554$, $p=3.4\times10^{-10}$）
    \item Energy：114.70M $\rightarrow$ 111.81M（$\Delta=-2.89M$, $p=0.086$）
    \item DSR：0.2734 $\rightarrow$ 0.2879（$\Delta=+0.0145$, $p=0.084$）
\end{itemize}

上层 attention-MAPPO 在 reward 和公平性维度上显著优于 uncoordinated greedy；energy 和 DSR 呈改善趋势，但未达到统计显著。

\textbf{补充表 A. 上层 vanilla MAPPO 消融（training seed=42, 4 workload seeds $\times$ 8 episodes）}

\begin{table}[!htbp]
    \centering
    \bicaption{上层 vanilla MAPPO 消融结果}{Upper-layer vanilla MAPPO ablation results}
    \label{tab:vanilla-ablation}
    \resizebox{\textwidth}{!}{
    \begin{tabular}{l|rr|rr|rr|rr}
        \hline
        策略 & Reward & & Latency & & Energy (M) & & DSR & \\
        \hline
        vanilla\_mappo\_\_heuristic & \textbf{-1462.6} & $\pm$ 328.7 & \textbf{841994} & $\pm$ 95239 & 120.11 & $\pm$ 15.69 & \textbf{0.2822} & $\pm$ 0.0278 \\
        attention\_mappo\_\_heuristic & -2222.1 & $\pm$ 890.6 & 1135453 & $\pm$ 132493 & \textbf{114.55} & $\pm$ 13.26 & 0.2799 & $\pm$ 0.0317 \\
        \hline
    \end{tabular}}
\end{table}

在该补充实验中，vanilla\_mappo 在 reward、latency 和 fairness 上优于 attention\_mappo；attention 相对 vanilla 的 latency 增加约 293458（paired t-test $p=0.0085$），reward 降低约 759.5（$p=0.0889$）。另一方面，attention-MAPPO 的 energy 更低（约 -5.57M，$p=0.102$），offline rate 更低，MBS load ratio 也显著更低（0.99\% vs 1.41\%，$p=0.0155$）。

\section{轨迹可视化分析}
\label{sec:trajectory-visualization}

由于本文方法同时包含上层轨迹控制和下层请求级卸载决策，仅依赖表格指标难以直观展示不同策略在空间维度上的行为差异。为此，图 \ref{fig:trajectory-user-distribution} 给出了同一评估场景下不同算法的无人机二维飞行轨迹。该图基于 training seed 42 和 workload seed 42 的第一个评估 episode 绘制。图中浅灰色散点表示评估 episode 初始时刻的 UE 分布，彩色曲线表示 5 架无人机的水平飞行轨迹，空心圆表示无人机起点，实心方块表示终点，箭头表示移动方向，虚线圆表示无人机终点处的覆盖范围，黑色五角星表示宏基站位置。

\begin{figure}[!htbp]
    \centering
    \begin{minipage}{0.48\textwidth}
        \centering
        \includegraphics[width=\linewidth]{docs/figures/图5-6a_随机策略无人机二维轨迹与用户分布.pdf}
        \\(a) 随机策略
    \end{minipage}
    \hfill
    \begin{minipage}{0.48\textwidth}
        \centering
        \includegraphics[width=\linewidth]{docs/figures/图5-6b_IPPO无人机二维轨迹与用户分布.pdf}
        \\(b) IPPO
    \end{minipage}

    \vspace{0.8em}

    \begin{minipage}{0.48\textwidth}
        \centering
        \includegraphics[width=\linewidth]{docs/figures/图5-6c_VanillaMAPPO无人机二维轨迹与用户分布.pdf}
        \\(c) Vanilla MAPPO
    \end{minipage}
    \hfill
    \begin{minipage}{0.48\textwidth}
        \centering
        \includegraphics[width=\linewidth]{docs/figures/图5-6d_本文方法无人机二维轨迹与用户分布.pdf}
        \\(d) 本文方法
    \end{minipage}
    \bicaption{不同算法下无人机二维轨迹与用户分布对比}{Two-dimensional UAV trajectories and UE distribution under different algorithms}
    \label{fig:trajectory-user-distribution}
\end{figure}

从图 \ref{fig:trajectory-user-distribution} 可以看出，随机策略的移动方向缺乏稳定目标，轨迹更多体现为无规则探索；IPPO 和 Vanilla MAPPO 在相同用户分布下形成了较明确的局部移动模式，其中 IPPO 更偏向快速靠近部分服务区域，因此在平均时延指标上表现较优。本文方法的轨迹并不以单纯缩短路径或最低时延为唯一目标，而是在用户分布、无人机间协作关系和下层卸载决策共同作用下形成相对稳定的空间部署。结合表 \ref{tab:main-results} 以及后续卸载分布分析可知，本文方法的优势主要体现在能耗、有效服务完成和协作卸载之间的综合权衡，而不是单一轨迹长度或单一时延指标最优。

\section{卸载分布分析}
\label{sec:offloading-analysis}

\begin{figure}[!htbp]
    \centering
    \includegraphics[width=0.85\textwidth]{docs/figures/fig_offloading_distribution.png}
    \bicaption{卸载分布对比}{Offloading distribution comparison}
    \label{fig:offloading-dist}
\end{figure}

从表 \ref{tab:main-results-cont} 的卸载比例数据可以看到：
\begin{itemize}
    \item Heuristic 基线（前两行）主要依赖 local（47\%）和 cooperative（50\%），MBS 使用极少（2.6-2.7\%）
    \item 下层 MAPPO 系列（后三行）显著改变了卸载分布：local 调整至 42-53\%，cooperative 调整至 30-45\%，MBS 增至 11-17\%
    \item Full hierarchical 实现了最均衡的卸载分布（local 46.5\%、cooperative 42.4\%、MBS 11.1\%），同时 offline rate 降至 0.00\%
\end{itemize}

\section{消融实验结果}
\label{sec:ablation-results}

表 \ref{tab:ablation-results} 给出了下层消融实验的详细结果。

\begin{table}[!htbp]
    \centering
    \bicaption{下层消融实验结果（相对于 lower\_full）}{Lower-layer ablation results (relative to lower\_full)}
    \label{tab:ablation-results}
    \begin{tabular}{l|rr|rr|rr}
        \hline
        消融 & Reward & & Energy (M) & & DSR & \\
        \hline
        lower\_full & -1599.4 & & 53.32 & & 0.2170 & \\
        no\_mask & -1556.3 & & 48.26 & & 0.1945 & \\
        no\_lagrange & -1585.9 & & 51.12 & & 0.1981 & \\
        no\_attention & -1626.7 & & 54.01 & & \textbf{0.2094} & \\
        \hline
    \end{tabular}
\end{table}

\textbf{质量感知 Mask（lower\_full vs no\_mask）}：
移除 mask 后 MBS load ratio 从 13.7\% 升至 18.0\%，策略变得对 MBS 更激进。Reward 略有上升（+43.1, $p=0.09$），但差异不显著。

\textbf{Lagrange 约束（lower\_full vs no\_lagrange）}：
移除 Lagrange 后 MBS 卸载比例从 13.7\% 飙升至 25.0\%（$\Delta=+11.3$pp, $p=6.5\times10^{-5}$）；对应地，MBS 负载比例从 6.3\% 升至 11.6\%。这说明 Lagrange 约束对 MBS 卸载依赖具有显著抑制作用。

\textbf{Attention 机制（lower\_full vs no\_attention）}：
移除 request attention 后 MBS load ratio 从 13.7\% 升至 17.3\%。DSR 为 0.2094，与 lower\_full 的 0.2170 接近。移除 attention 后 cooperative ratio 从 44.6\% 降至 34.4\%，local ratio 从 41.7\% 升至 48.2\%。attention 机制使策略能更灵活地利用 cooperative 卸载机会。

\begin{figure}[!htbp]
    \centering
    \includegraphics[width=0.85\textwidth]{docs/figures/fig_ablation_comparison.png}
    \bicaption{消融实验对比}{Ablation experiment comparison}
    \label{fig:ablation-compare}
\end{figure}

\section{端到端 joint MAPPO baseline}
\label{sec:joint-mappo}

为回应"为什么不直接训练单一端到端策略同时输出轨迹和卸载动作"的问题，本文实现了 \texttt{joint\_mappo} baseline。该模型使用单一 MAPPO actor 同时输出 UAV 连续轨迹动作和请求级离散卸载动作，并使用联合观测与统一 PPO 更新，不复用已训练的上层或下层模型。

\textbf{补充表 B. 端到端 joint MAPPO 评估结果（training seed=42, 4 workload seeds $\times$ 8 episodes）}

\begin{table}[!htbp]
    \centering
    \bicaption{端到端 joint MAPPO 评估结果}{End-to-end joint MAPPO evaluation results}
    \label{tab:joint-mappo}
    \begin{tabular}{l|rr|rr|rr|rr}
        \hline
        策略 & Reward & & Latency & & Energy (M) & & DSR & \\
        \hline
        joint\_end\_to\_end\_mappo & -9865.1 & $\pm$ 325.6 & 1580604 & $\pm$ 267198 & 81.58 & $\pm$ 31.25 & 0.1349 & $\pm$ 0.0761 \\
        \hline
    \end{tabular}
\end{table}

训练阶段 \texttt{joint\_mappo} 在 200 episodes 后的 recent reward 为 -4030.8，但独立 workload 评估中的平均 reward 为 -9865.1，明显低于分层策略；同时 DSR 仅为 0.1349，fairness 为 0.4581。该结果表明，在相同训练预算下，直接端到端联合学习面临更大的联合动作空间、混合连续-离散动作建模和 credit assignment 难度，样本效率和稳定性均弱于双层分解方案。

\section{DSR-priority 主配置实验}
\label{sec:dsr-priority}

为进一步验证框架在服务质量偏好上的可调节性，本文将 DSR-priority 作为第二组正式主配置进行评估。相较能耗优先配置，DSR-priority 配置将系统级权重调整为 $W_{DSR}=3.0$、$W_E=0.2$；下层 Lagrange 约束目标调整为 $\tau_{dsr}=0.28$、$\tau_{mbs}=0.08$，Lagrange 学习率为 $\eta_{\lambda}=0.5$，乘子上界为 $\lambda_{max}=50.0$；下层 reward 权重调整为 $w_{dead}=7.0$, $w_{succ}=1.5$, $w_{lat}=2.5$, $w_{en}=0.5$, $w_{mbs}=0.2$。

\begin{table}[!htbp]
    \centering
    \bicaption{三种配置下的 Full Hierarchical MARL 对比（$N=30$）}{Full Hierarchical MARL under three configurations ($N=30$)}
    \label{tab:dsr-priority}
    \begin{tabular}{l|rr|rr|rr}
        \hline
        配置 & Reward & & Energy (M) & & DSR & \\
        \hline
        baseline（heuristic） & -5933.9 & & 114.70 & & 0.2734 & \\
        能耗优先（base v2） & -1451.7 & & \textbf{58.06} & & 0.2083 & \\
        DSR 优先（strong） & \textbf{-1147.2} & & 83.04 & & \textbf{0.2435} & \\
        \hline
    \end{tabular}
\end{table}

DSR-priority 配置下，完整方法的 DSR 从能耗优先配置下的 0.2083 提升至 0.2435，与 baseline 的差距从 -0.065 缩小至 -0.030（$p = 1.6\times10^{-6}$）。但该值仍低于 heuristic baseline 的 0.2734，因此本文将其解释为 DSR 的部分恢复，而非完全恢复。DSR-priority 配置下能耗为 83.04M，相比 baseline 的 114.70M 仍节省 27.6\%；cooperative ratio 达到 56.7\%，为所有学习式策略中最高。

\begin{figure}[!htbp]
    \centering
    \includegraphics[width=0.85\textwidth]{docs/figures/fig_lagrange_dsr_mbs_tradeoff.png}
    \bicaption{Lagrange 约束动态与 DSR-MBS 权衡}{Lagrange constraint dynamics and DSR-MBS trade-off}
    \label{fig:lagrange-tradeoff}
\end{figure}

\begin{figure}[!htbp]
    \centering
    \includegraphics[width=0.85\textwidth]{docs/figures/fig_hierarchical_training_curves.png}
    \bicaption{固定 workload42 评估曲线}{Fixed workload42 evaluation curves}
    \label{fig:training-curves}
\end{figure}

\section{讨论}
\label{sec:discussion}

\subsection{Energy-DSR-Fairness 三元权衡}
\label{subsec:tradeoff}

实验揭示了多 UAV MEC 系统中一个核心的三元权衡：

\begin{table}[!htbp]
    \centering
    \bicaption{Energy-DSR-Fairness 三元权衡分析}{Energy-DSR-Fairness ternary trade-off analysis}
    \label{tab:tradeoff}
    \begin{tabular}{|l|l|l|}
        \hline
        优化方向 & 优势 & 代价 \\
        \hline
        降低 Energy & 从 114.7M $\rightarrow$ 42.8M (-63\%) & DSR 从 0.27 $\rightarrow$ 0.19 \\
        提升 DSR & 0.27 $\rightarrow$ 0.29 & Energy 维持高位 111.8M \\
        提升 Fairness & 0.77 $\rightarrow$ 0.93 & 依赖更强的轨迹协同建模 \\
        \hline
    \end{tabular}
\end{table}

本文的双层 MARL 框架通过 Lagrange 约束和奖励权重提供了调控这一权衡的机制。第 \ref{sec:dsr-priority} 节的 DSR-priority 主配置实验表明：将 $W_{DSR}$ 从 1.0 提升至 3.0、$W_E$ 从 0.5 降至 0.2、$\tau_{dsr}$ 从 0.18 提升至 0.28 后，完整方法的 DSR 可从 0.2083 部分恢复至 0.2435，与 baseline 的差距缩小约 54\%，同时仍保持约 28\% 的能耗节省。

\subsection{上层轨迹 vs 下层卸载的相对贡献}
\label{subsec:layer-contribution}

通过对比 \texttt{att+heuristic}（仅上层）和 \texttt{unco+lower}（仅下层）可以发现：
\begin{itemize}
    \item 上层贡献主要在于 \textbf{公平性}（0.930 vs 0.775）和 \textbf{整体奖励}（-2129 vs -5934）
    \item 下层贡献主要在于 \textbf{能耗降低}（42.9M vs 114.7M）和 \textbf{卸载灵活性}
    \item 完整双层（-1451.7, 58.06M, 0.936, 0.208）综合了两层优势
\end{itemize}

\subsection{与现有工作的对比}
\label{subsec:comparison}

相比黄子祥等 \cite{huang2026madr} 基于 MADRL 的单层端到端方法，本文的双层分解将轨迹控制与请求级卸载拆开处理，降低了直接联合动作建模的结构复杂度，并提高了策略行为的可解释性。相比尤昕阳等 \cite{you2026soft} 的 SAC 方法，本文通过双层拆分分别处理连续轨迹控制与离散请求卸载，减少了直接构造混合联合动作的难度。相比曾耀平等 \cite{zeng2026uav} 的博弈论方法，本文基于 MARL 的在线执行主要依赖局部和邻居观测，具有较好的动态适应性。

\section{本章小结}
\label{sec:summary-chap05}

本章介绍了实验设计和实验结果分析。首先介绍了仿真环境、对比方法和评估指标，然后通过主实验、上层贡献分析、卸载分布分析、消融实验、端到端 baseline 实验和 DSR-priority 配置实验，全面验证了所提双层 MARL 框架的有效性。实验结果表明，该框架在奖励、能耗、公平性和离线率等指标上均取得了显著改善，并通过 Lagrange 约束和奖励权重提供了调控 DSR-能耗权衡的机制。
